\chapter{Deforming Kleinian manifolds by homeomorphisms of the sphere at infinity}

\section{Extensions of vector fields}

\begin{definition}[Visual averaging of boundary vector fields]
\label{thurston-ch11-visual-average}
\group{extensions}
\source{thurston:page-0302,thurston:page-0303}
\notready
For $x\in H^3$, identify the sphere at infinity with the unit sphere in
$T_xH^3$.  For $y\in S^2_\infty$, let
$i_{x,y}:T_yS^2_\infty\to T_xH^3$ be the induced identification.  The visual
average of a boundary vector field $X$ is
\[
 \operatorname{av}X(x)=\frac{1}{4\pi}
 \int_{S^2_\infty}i_{x,y}(X(y))\,dV_x(y),
\]
where $V_x$ is visual measure at $x$.
\end{definition}
\begin{proof}
\sketch The visual identification transports every value of $X$ into the
single vector space $T_xH^3$, so the normalized vector-valued integral is
well-defined.
\end{proof}

\begin{lemma}[Upper-half-space visual measure]
\label{thurston-ch11-visual-measure}
\group{extensions}
\source{thurston:page-0303}
\notready
In the upper half-space model, if $x=(z_0,h)$ and a geodesic from $x$ meets
$\mathbb C$ at radial distance $r$, then
\[
 r=h\cot(\theta/2),\qquad
 dV_x=4\left(h+\frac{r^2}{h}\right)^{-2}d\mu.
\]
The boundary field induced by a tangent vector $X$ at $x$ is the restriction
of the infinitesimal hyperbolic translation in direction $X$.
\end{lemma}
\begin{proof}
\sketch Differentiate $r=h\cot(\theta/2)$ to obtain
$dr=-\frac12(h+r^2/h)d\theta$.  Conformality of the visual map then gives
the displayed density.  The translation description follows either from the
disk model or, for a vertical unit vector, from
$-\sin\theta\,\partial_\theta=r\,\partial_r$.
\end{proof}

\begin{proposition}[Renormalized extension]
\label{thurston-ch11-renormalized-extension}
\dcref{11.1.1}
\group{extensions}
\source{thurston:page-0304,thurston:page-0305,thurston:page-0306}
\notready
Define
\[
 \operatorname{ex}(X)=\frac32\operatorname{av}(X)\quad\text{on }H^3,
 \qquad \operatorname{ex}(X)=X\quad\text{on }S^2_\infty.
\]
If $X$ is continuous or Lipschitz on $S^2_\infty$, then $\operatorname{ex}(X)$
has the same regularity on the closed ball.  If $X$ is holomorphic, then
$\operatorname{ex}(X)$ is the infinitesimal hyperbolic isometry extending $X$.
\end{proposition}
\begin{proof}
\sketch For a horizontal unit vector the visual integral is
$\operatorname{av}(\partial_z)=(2/3h)\partial_z$, so the factor $3/2$
gives the correct boundary scale.  Parabolic fields extend correctly, and
the fields $\partial_z,z^2\partial_z,(z-1)^2\partial_z$ span all holomorphic
quadratic vector fields.  Translate a point tending to the boundary to the
origin; continuity gives approximation by a parabolic field.  A boundary
Lipschitz bound gives a uniform bound after this translation, and scaling back
shows a uniform Lipschitz bound both at the boundary and in the interior.
\end{proof}

\begin{definition}[Quasi-isometric vector field]
\label{thurston-ch11-quasi-isometric-field}
\group{extensions}
\source{thurston:page-0307}
\notready
A vector field is $k$-quasi-isometric if its flow $\varphi_t$ satisfies
\[
 e^{-kt}d(x,y)\le d(\varphi_t(x),\varphi_t(y))\le e^{kt}d(x,y).
\]
\end{definition}
\begin{proposition}[Boundary isotopies extend quasi-isometrically]
\label{thurston-ch11-isotopy-extension}
\dcref{11.1}
\group{extensions}
\source{thurston:page-0307}
\notready
A boundary isotopy of a Kleinian manifold whose time derivative is Lipschitz
extends canonically by $\operatorname{ex}$ to the interior.  Its extension has
quasi-isometric time derivative; in particular, the time-one map is a
quasi-isometry.  A $k$-Lipschitz vector field is $k$-quasi-isometric.
\end{proposition}
\begin{proof}
\sketch Lipschitz vector fields have unique flows, and the local proof of the
extension proposition applies on quotient charts.  The differential inequality
for distances follows from the Lipschitz bound; translating an arbitrary point
to the origin and comparing with a parabolic field gives the same estimate
throughout hyperbolic space.
\end{proof}

\begin{proposition}[Derivative of the extension]
\label{thurston-ch11-extension-derivative}
\dcref{11.1.2}
\group{derivatives}
\source{thurston:page-0307}
\notready
The extension is natural under hyperbolic isometries:
$\operatorname{ex}(T_*X)=T_*\operatorname{ex}(X)$.  If $Y$ is an infinitesimal
translation whose axis passes through $x$, then
\[
 \nabla_{Y_x}\operatorname{ex}X=\operatorname{ex}[Y,X].
\]
\end{proposition}
\begin{proof}
\sketch Differentiate naturality at the identity to obtain
$\operatorname{ex}[Y,X]=[Y,\operatorname{ex}X]$.  Along the axis,
$\nabla_XY=0$, and the identity
$[Y,X]=\nabla_YX-\nabla_XY$ yields the formula.
\end{proof}

\begin{definition}[Derivative decompositions]
\label{thurston-ch11-derivative-decompositions}
\group{derivatives}
\source{thurston:page-0308,thurston:page-0309,thurston:page-0310}
\notready
For a vector field $X$ on $H^3$, write
$\nabla X=\nabla^sX+\nabla^aX$, where
\[
 \nabla^s_YX\cdot Y'=\tfrac12(\nabla_YX\cdot Y'+\nabla_{Y'}X\cdot Y),\quad
 \nabla^a_YX\cdot Y'=\tfrac12(\nabla_YX\cdot Y'-\nabla_{Y'}X\cdot Y).
\]
Define $\operatorname{curl}X$ by
$\nabla^a_YX=\frac12\operatorname{curl}X\times Y$, and decompose
$\nabla^sX=\frac13(\operatorname{div}X)I+\nabla^{s0}X$.
For a surface vector field define
\[
 \partial X(Y)=\tfrac12(\nabla_YX-i\nabla_{iY}X),\qquad
 \bar\partial X(Y)=\tfrac12(\nabla_YX+i\nabla_{iY}X).
\]
\end{definition}
\begin{proof}
\sketch The first decomposition is the symmetric/antisymmetric decomposition
of a linear map, while the trace-free split removes the scalar part.  Splitting
a real tangent map into its complex-linear and conjugate-linear parts gives
the definitions of $\partial$ and $\bar\partial$.
\end{proof}

\begin{proposition}[Curl identities for visual extensions]
\label{thurston-ch11-extension-curl}
\dcref{11.1.3}
\group{derivatives}
\source{thurston:page-0308,thurston:page-0309}
\notready
For every boundary vector field $X$,
\[
 \operatorname{curl}(\operatorname{ex}X)=2\operatorname{ex}(iX),\qquad
 \operatorname{curl}^2(\operatorname{ex}X)=-4\operatorname{ex}X,\qquad
 \operatorname{div}(\operatorname{ex}X)=0.
\]
\end{proposition}
\begin{proof}
\sketch In horospherical coordinates, an infinitesimal visual contribution
has the form $Z=e^{-3t}e_1$ and the connection gives
$\operatorname{curl}Z=-2e^{-3t}e_2$.  Integrating contributions identifies the
first formula, the second follows by applying it to $iX$, and the third from
$\operatorname{div}\operatorname{curl}=0$.
\end{proof}

\begin{proposition}[Complex derivative decomposition]
\label{thurston-ch11-complex-derivative}
\dcref{11.1.4}
\group{derivatives}
\source{thurston:page-0310,thurston:page-0311}
\notready
On an oriented surface,
\[
 \partial X=\tfrac12(\operatorname{div}X)I+\nabla^aX
 =\tfrac12\{(\operatorname{div}X)I+(\operatorname{curl}X)iI\},
 \qquad \bar\partial X=\nabla^{s0}X.
\]
The operator $\partial X$ is invariant under conformal changes of metric.
\end{proposition}
\begin{proof}
\sketch Any real linear map is the sum of its complex-linear and conjugate-
linear parts, and comparison with the trace, antisymmetric, and trace-free
components gives the two identities.  Using torsion-freeness,
$\bar\partial X(Y)=\tfrac12([Y,X]+i[iY,X])$ removes the metric connection and
proves conformal invariance.
\end{proof}

\begin{proposition}[Symmetric derivative of an extension]
\label{thurston-ch11-extension-symmetric-derivative}
\dcref{11.1.5}
\group{derivatives}
\source{thurston:page-0311,thurston:page-0312}
\notready
For $Y\in T_xH^3$ and a $C^1$ boundary vector field $X$,
\[
 \nabla^s_Y\operatorname{ex}X
 =\frac{3}{4\pi}\int_{S^2_\infty}i_*\bigl(\bar\partial X(Y_\infty)\bigr)\,dV_x.
\]
\end{proposition}
\begin{proof}
\sketch Both sides are symmetric in the tested tangent direction.  Apply the
derivative formula to $Y$ and use that the extension of $iY_\infty$ vanishes
at $x$; the resulting two integral identities combine into the displayed
integral by the definition of $\bar\partial$.
\end{proof}
